\section{Discussion}
\label{sec:discussion}

\paragraph{Baseline choice.}
Ensemble MC dropout provides a consistently strong neural reference, but its best average rank does
not establish an advantage over the other ensemble pipelines. LightGBM's many individual wins and
poor results elsewhere show why consistency across a benchmark and performance on a particular
assay can differ. The per-dataset comparisons identify where one method has a clear advantage and
where several alternatives cannot be ruled out as best. A comparison using only one baseline family
can therefore give an incomplete picture of a new method's performance. Both LightGBM and an
ensemble neural baseline should be included when evaluating new molecular predictors, including
pretrained models.

\paragraph{Variation across evaluations.}
Evaluation introduces uncertainty of its own. Changing which compounds enter training and testing
can change the apparent leader, and averaging model seeds on one fold leaves variation between
chemical partitions unexplored. At the same number of fold-score evaluations, completing
cross-validation gives a more stable full rank ordering than repeating seeds on one fold. A compact
evaluation can help screen methods; a claim of improvement should be checked across multiple
chemical partitions. Paired comparisons account for the variation in method differences, while
multiplicity adjustment limits spurious claims of a difference when several methods are compared.

Shared training compounds weaken the independence assumption behind the per-dataset tests. Some
endpoint families also share compounds, and only nine datasets contribute to the cross-dataset
interval-score comparison. Its unresolved differences deserve caution: the absence of a detected
difference gives little reason to prefer one method over another on statistical grounds alone.

\paragraph{Tuning and computational cost.}
Validation data affect the comparison before test scores are calculated. A single validation split
often selects a different configuration from five-fold selection, and those choices perform worse
across the five validation folds. Authors proposing a new model should give simple baselines a fair
chance through careful hyperparameter and stopping-point selection. Their competitiveness here makes
that tuning part of establishing an improvement, including for expensive pretrained architectures.

Ensembles require additional model fits, while other UQ methods add posterior sampling or last-layer
optimization. We made the full benchmark feasible by sharing MLP hyperparameter and checkpoint
selection across UQ methods. These choices optimize point prediction. Selecting the final
ensemble or using a calibrated validation score could favor different hyperparameters or stopping
epochs \citep{ranjan2024posthoc,fredsgaard2025joint}. The rankings consequently describe the complete
pipelines evaluated here, including their selection procedures and additional computation.

\paragraph{Using the uncertainty estimates.}
The relevant score depends on what the predictions will be used for. Brier score assesses
probabilities through both calibration and refinement without separating epistemic and aleatoric
uncertainty. The regression results depend on the calibration maps and Gaussian assumptions used to
construct the intervals. An application that uses uncertainty mainly to rank candidates would also
need to assess that ranking.

Claims of improvement should be supported by comparisons with both carefully tuned baselines.
Reporting the size of the gain and its variation across chemical partitions helps readers judge
its practical value.
